\section{Gêmeos Digitais na Saúde Pública (SUS)}

Em um cenário demográfico continental como o Brasil, a verdadeira revolução do Gêmeo Digital transcende a arcada individual e atinge a malha epidemiológica. 

\subsection{O Gêmeo Digital Populacional (Epidemiologia \textit{in silico})}

Enquanto o gêmeo digital clínico modela o indivíduo, o \textbf{Gêmeo Digital Populacional} modela uma comunidade inteira. Ao integrar dados do inquérito SB Brasil, do Sistema de Informação Ambulatorial (SIA-SUS) e índices de saneamento básico, o Ministério da Saúde poderá simular a eficácia de intervenções macroestruturais. Por exemplo, antes de aprovar uma verba milionária para centros de especialidade, o algoritmo simula em ambiente virtual se o maior retorno financeiro e social estaria na contratação de protesistas ou no investimento em fluoretação de águas na região analisada \cite{frota2025sus}.

\subsection{Escalando a Tele-Odontologia de Precisão}

Para as Unidades de Saúde da Família (USF), a tele-odontologia equipada com modelagem virtual reduzirá a fila de espera para especialidades críticas (Endodontia e Cirurgia Bucomaxilofacial). A democratização do escaneamento permite que a atenção básica seja o pólo de captação digital, transformando cirurgiões-dentistas generalistas em executores de planos de voo milimetricamente traçados por IAs em grandes centros de referência nacionais \cite{d24am2026gemeos}.
